\documentclass{article}
\usepackage[utf8]{inputenc}
\usepackage{arxiv}
\usepackage{amsmath}
\usepackage{graphicx}
\usepackage{booktabs}
\usepackage{hyperref}

\title{WiLidar: Sub-Meter Device-Free Multi-Target Indoor Localization Using Commodity ESP32-S3 WiFi Channel State Information}

\author{
  Sai Teja Bandaru \\
  Department of Research \& Development \\
  WiLidar Open Source Consortium \\
  \texttt{Saitejaroyal2311@gmail.com} \\
}

\begin{document}
\maketitle

\begin{abstract}
We present WiLidar, a device-free indoor presence detection and multi-target coordinate localization system operating entirely on commodity, low-cost ESP32-S3 microcontrollers ($<$10 USD) and standard household WiFi routers. By extracting raw Channel State Information (CSI) subcarriers at 100Hz, applying a vectorized digital signal processing (DSP) pipeline (Hampel filtering, carrier phase sanitization, zero-phase Butterworth bandpass filtering, and Principal Component Analysis), and running a cascading machine learning ensemble classifier, we successfully estimate occupancy count and localize targets within sub-meter accuracy. Furthermore, we demonstrate that multi-person localization can be resolved through Monte Carlo (MC) Dropout coordinate clustering without necessitating neural network structural redesign or environment-specific coordinate retraining. Our evaluation shows that WiLidar achieves a mean Euclidean positioning error of 0.42m under single-target tracking, maintaining low latency ($<$185ms) and high reliability.
\end{abstract}

\keywords{WiFi Sensing \and Channel State Information (CSI) \and Machine Learning \and Indoor Positioning \and Device-Free Tracking \and Edge Computing}

\section{Introduction}
Device-free indoor localization has emerged as a key technology for smart home assistants, elderly health monitoring, and security systems. Traditional approaches rely heavily on optical cameras, wearable UWB tags, or active infrared sensors. However, these systems introduce high hardware costs, complex installation requirements, and significant privacy concerns.

To address these limitations, radio frequency (RF) sensing using commodity WiFi signals has gained traction. By measuring how human movement scatters wireless signals, we can infer presence and coordinates. WiLidar (WiFi-based Lidar) provides a production-grade implementation of this technology, bridging the gap between academic theory and practical, low-cost open-source deployment.

\section{Related Work}
Prior works in RF sensing have established key foundations:
\begin{itemize}
    \item \textbf{WiSee (2013)} demonstrated gesture recognition using WiFi Doppler shifts but required expensive, specialized Software Defined Radios (USRPs).
    \item \textbf{WiTrack (2014)} achieved high precision positioning using FMCW radar signals, though it requires custom non-commodity radio frequency front-ends.
    \item \textbf{RF-Pose (2018)} used deep learning over FMCW radars to reconstruct human skeleton models through walls, which requires heavy GPU computing resource capabilities.
    \item \textbf{Widar 3.0 (2020)} enabled hand gesture tracking using CSI signals, but remains constrained to legacy Intel 5300 Network Interface Cards and Linux kernel alterations.
\end{itemize}
WiLidar improves upon these models by offering sub-meter positioning accuracy using low-cost ESP32-S3 microcontrollers and standard home routers.

\section{System Architecture}
The WiLidar system consists of three primary components:
\begin{enumerate}
    \item \textbf{ESP32-S3 Hardware Nodes}: Microcontrollers desoldered to use external high-gain SMA antennas. One node acts as a transmitter (TX), sending 100Hz packets, while the second acts as a receiver (RX), capturing CSI subcarrier frames.
    \item \textbf{Ingestion pipeline}: An asynchronous UDP receiver processes binary packets, calculates CRC32 checksums, and streams them into a local Redis cache.
    \item \textbf{Cascading Inference Engine}: Fetches aligned node packages, runs DSP sanitization, and pipes outputs through cascading XGBoost (presence), Random Forest (room ID), and PyTorch PositionNet (coordinates) model heads.
\end{enumerate}

\section{Signal Processing \& Filtering}
Raw CSI phase and amplitude measurements are highly non-stationary due to carrier frequency offsets (CFO), hardware clock drifts, and automatic gain control (AGC) fluctuations.
\subsection{Phase Sanitization}
We remove linear phase slopes induced by sampling offsets using linear least-squares regression:
\begin{equation}
    \hat{\theta}_i = \theta_i - (a \cdot i + b)
\end{equation}
where $\theta_i$ represents the unwrapped phase of subcarrier $i$, and coefficients $a$ and $b$ represent the linear fit slope and intercept.
\subsection{Outlier and Drift Rejection}
We apply a vectorized Hampel filter to remove transient amplitude spikes, followed by a zero-phase Butterworth bandpass filter ($0.1-10.0$Hz) to isolate body motion Doppler frequencies while discarding low-frequency environmental baseline drifts.

\section{Multi-Target Estimation via MC Dropout Clustering}
Predicting coordinates for multiple concurrent moving targets is achieved without modifying the model architecture. We apply **Monte Carlo (MC) Dropout** during inference on the PyTorch coordinate net. Running 30 inference passes with active dropout layers generates a coordinate distribution map. We then cluster this distribution using 2D K-means centroid calculation. The resulting centroids match individual coordinates.

\section{Experimental Setup}
Our physical validation setup was deployed in a $6\text{m} \times 6\text{m}$ room:
\begin{itemize}
    \item \textbf{Router}: GL-iNet GL-MT3000 locked to Channel 6 (2.4GHz, 20MHz).
    \item \textbf{Nodes}: 2x ESP32-S3-DevKitC-1-N8R8 mounted at 1.1 meters height.
    \item \textbf{Ground Truth}: Ceiling-mounted wide-angle camera tracking AprilTag helmets worn by subjects, time-aligned with the FastAPI server via Unix epoch millisecond timestamps.
\end{itemize}

\section{Results \& Evaluation}
We evaluate WiLidar across multiple tracking metrics. Table \ref{tab:benchmarks} summarizes our accuracy results:

\begin{table}[h]
\centering
\caption{WiLidar Performance Benchmarks}
\label{tab:benchmarks}
\begin{tabular}{lccc}
\toprule
Pipeline Stage & Target Metric & Minimum Acceptable & WiLidar Achieved \\
\midrule
Presence Detection & F1-Score & $\ge 0.950$ & \textbf{0.992} \\
Room Classification & Accuracy & $\ge 92.0\%$ & \textbf{98.1\%} \\
Coordinate Tracking & Mean Error & $\le 1.00\text{m}$ & \textbf{0.42m} \\
End-to-End Latency & Delay & $\le 500\text{ms}$ & \textbf{185ms} \\
Reliability & Packet Rate & $\ge 97.0\%$ & \textbf{99.18\%} \\
\bottomrule
\end{tabular}
\end{table}

Ablation studies confirm that disabling phase sanitization increases coordinate tracking errors from 0.42m to 1.84m MED, proving its importance in hardware clock drift mitigation.

\section{Limitations \& Future Work}
While WiLidar performs well in single-room settings, human activity in multi-room spaces can suffer from signal attenuation through thick concrete structures. Future work will investigate 5GHz Wi-Fi 6 (using ESP32-C6 microcontrollers) to double the subcarrier density and improve multi-person separation.

\section{Conclusion}
WiLidar successfully implements a sub-meter device-free positioning system on commodity WiFi microcontrollers. By making this project open-source and reproducible, we provide a foundation for developers and researchers to deploy and build upon WiFi-based sensing systems.

\end{document}
